%!TEX root = ../../thesis.tex

\section{The Design Space}\label{sec:mdp-design-space}

% TODO: Lay out the space, then place the two formulations in it. This is the map the
% reader carries into \cref{sec:input-actions,sec:corpus-actions}, so use the same
% axis names as the subsection headings there:
%   - action granularity: one input symbol per step vs.\ one corpus operation per step,
%   - what the observation must expose (raw bytes / coverage / per-slot features),
%   - reward: absolute per-execution coverage vs.\ marginal new coverage,
%   - episode structure: fixed input length vs.\ bounded iteration budget,
%   - how determinism is obtained (for free vs.\ bought with an exhaustive sweep),
%   - executions per agent decision (1 vs.\ up to 256) --- the lever that decides
%     whether MCTS overhead is amortized.
%
% Then state both formulations side by side: one table, one row per axis above, one
% column per formulation. That table is the spine of this chapter --- the two
% following sections are its columns expanded, in the same order, and
% \cref{sec:results-head-to-head} compares them empirically on one target.
%
% Close with the hypothesis the comparison tests (RQ2, \cref{sec:introduction-rqs})
% and with the lineage of each column: input-level descends from the protocol and
% state-selection line \cite{borcherding_bandits_2023}, corpus-level from
% \cite{bottinger_deep_2018}. Keep the lineage to one sentence each --- the full
% comparison is \cref{sec:mdp-prior-work}.
